\documentclass[11pt]{article}

\usepackage[margin=1in]{geometry}
\usepackage{enumitem}
\usepackage{amsmath}
\usepackage{hyperref}
\usepackage{xcolor}
\usepackage{titlesec}
\usepackage{parskip}

\hypersetup{
    colorlinks=true,
    linkcolor=blue,
    urlcolor=blue
}

\titleformat{\section}{\large\bfseries}{\thesection}{1em}{}
\titleformat{\subsection}{\normalsize\bfseries}{\thesubsection}{1em}{}

\title{\textbf{Masting Project -- UBC Visit} \\[0.3em]
\large January--March 2026 --- Analysis Plan (Draft)}
\author{}
\date{}

\begin{document}

\maketitle
\vspace{-2em}

\noindent\textbf{People involved:} Victor, Janneke, Lizzie, and Él\'eonore \textit{(Am I forgetting someone? Is Mao involved?)}

\noindent\textbf{Study system:} Mount Rainier National Park

\noindent\textbf{Dataset:} Seed-fall data (multi-species, multi-site, multi-year)

\section*{Research Questions and Hypotheses}

\subsection*{1) Do species exhibit interspecific synchrony in seed production at Mount Rainier National Park?}

\textbf{Hypothesis:} Preliminary results suggest interspecific synchrony among several species.

\textbf{Approach:}
\begin{itemize}[leftmargin=1.5em]
    \item For each site, calculate yearly mean seed production per species.
    \item Quantify synchrony using the variance ratio (Loreau \& de Mazancourt) applied to species-level time series.
    \item Compare observed synchrony against null expectations?
\end{itemize}

\subsection*{2) Is seed production synchronized within species (intraspecific synchrony) across sites at Mount Rainier?}

\textbf{Approach:}
\begin{itemize}[leftmargin=1.5em]
    \item For each species, compare annual seed production across sites.
    \item Apply the variance ratio to multisite time series for that single species.
    \item Compare synchrony levels between species to see which species show tighter intraspecific coupling.
\end{itemize}

\subsection*{3) Is total resource availability (g/m\textsuperscript{2}) temporally synchronized from the perspective of seed predators?}

\textbf{Note:} Synchrony in total resource availability may drive predator responses and link to the seed predator satiation hypothesis.

\textbf{Approach:}
\begin{itemize}[leftmargin=1.5em]
    \item Define a metric for resource availability (e.g., annual total seed fall per m\textsuperscript{2} or a nutritional index if species differ strongly in energetic value).
    \item Identify years of low vs.\ high total availability (resource pulses).
    \item Explore how these pulses might relate conceptually to chapters 2 and 3 (predation and predator satiation), without going too far analytically yet.
\end{itemize}

\subsection*{4) (Potential) How does interspecific synchrony vary along the elevational gradient?}

\textbf{Approach:}
\begin{itemize}[leftmargin=1.5em]
    \item Estimate interspecific variance ratios separately for each elevation band (site).
    \item Examine whether synchrony strengthens or weakens with elevation.
\end{itemize}

\subsection*{5) (Potential) How does intraspecific synchrony vary with elevation, and can these patterns be linked to seed-predator home-range sizes?}

\textbf{Approach:}
\begin{itemize}[leftmargin=1.5em]
    \item Quantify intraspecific synchrony for each species across elevations.
    \item If differences emerge, explore whether predator movement ecology (e.g., home-range size of rodents, corvids, squirrels) could interact with the scale of resource pulses.
\end{itemize}

\section*{Key Points to Clarify}

\begin{itemize}[leftmargin=1.5em]
    \item \textbf{Terminology:} We are not defining masting here; instead, we treat species as masting or quasi-masting based on known ecology. We can either:
    \begin{itemize}
        \item use ``synchrony in seed production'' throughout, \textbf{OR}
        \item refer to ``masting synchrony'' without formally quantifying masting itself (unless CV is needed).
    \end{itemize}
    \item \textbf{Modeling:} Adapt Victor's model to handle multiple species jointly for interspecific analyses.
    \item \textbf{Resource availability metric:} Must be defined early---simple biomass (g/m\textsuperscript{2}) may be sufficient, but consider adjusting for seed energy content if relevant.
\end{itemize}

\end{document}
